\section{Conclusion}

Phase 2 is a direct continuation of the completed Phase 1 DOA-sector inference project. Phase 1 established that compact beamspace features from a physics-based ULA model can support both continuous nearest-centroid classification and rate-coded LIF-inspired temporal inference. The resulting $86.90\%$ continuous accuracy and $84.17\%$ best LIF accuracy demonstrated feasibility and exposed a meaningful accuracy--latency tradeoff, but those results were obtained primarily under nominal and aggregated sensing conditions.

The proposed Phase 2 work advances the project from feasibility to controlled robustness evaluation. It will preserve the Phase 1 signal model, covariance and beamspace pipeline, classifier prototypes, rate encoder, and LIF accumulator while introducing three explicit stress dimensions: extended low-SNR operation, sector-boundary proximity, and modest element-level gain and phase mismatch. Bartlett sector estimation, continuous nearest-centroid classification, and spike-based temporal inference will be compared on identical test samples using common sector-level metrics.

The principal engineering outcome will be a quantitative account of accuracy, confusion behavior, stable-decision latency, spike activity, and accuracy retention across nominal and stressed conditions. A renewed LIF parameter sweep will identify Pareto-relevant operating points that balance accuracy retention, temporal cost, and sparse activity. The integrated final run will then determine whether the selected operating point remains defensible when multiple sensing difficulties occur together.

Phase 2 does not claim a new high-resolution DOA estimator, a complete radar processor, biological SNN fidelity, direct energy savings, or a finished hardware implementation. Its contribution is the reproducible comparative evidence needed to decide whether the Phase 1 spike-based architecture is sufficiently robust and technically justified for a later transition to fixed-point analysis, PYNQ deployment, FPGA acceleration, or Vitis HLS implementation.

The completed phase will be delivered as a traceable package containing source code, named experiment runs, numerical outputs, report-ready figures, a formal technical report, and an indexed handoff document. This structure ensures that Phase 2 remains aligned with Phase 1 while creating a disciplined engineering foundation for the next embedded-AI implementation stage.